\section{Interfaces}
\lstset{ % Allgemeine Einstellungen für alle Codeblöcke
    language=C,% Sprache für Syntax-Highlighting
    basicstyle=\ttfamily, % Grundlegender Schriftstil
    keywordstyle=\color{blue},     % Farbe der Schlüsselwörter
    commentstyle=\color{gray},     % Farbe der Kommentare
    stringstyle=\color{red},       % Farbe der Zeichenketten
    numbers=left,                  % Zeilennummern links anzeigen
    numberstyle=\tiny\color{gray}, % Stil der Zeilennummern
    frame=single,                  % Rahmen um den Code
    breaklines=true,               % Automatischer Zeilenumbruch
    captionpos=b, % Bildunterschrift unterhalb des Codes
}

\begin{frame}[fragile] %fragile tag benötigt, da code Listings auftreten
\frametitle{Interfaces}
Was sind Interfaces?
\begin{itemize}
    \item Von Interfaces können keine Objekte erstellt werden
    \item Sie beschreiben die Schnittstelle einer Klasse, d.h. welche Methoden sie besitzen muss
    \item können von Klassen implementiert werden
    \item schreiben Klassen vor, welche Methoden sie besitzen müssen
    \item eine Klasse kann beliebig viele Interfaces implementieren
    \item Methoden in Interfaces sind standardmäßig public und abstract, d.h. sie müssen von der Klasse implementiert werden
\end{itemize}
Beispiel für ein Interface:
\begin{lstlisting}
    interface Item {
        double weight();
    }
\end{lstlisting}\footnote{\textit{public} ist in Methoden, die von einem Interface vorgeschrieben werden nicht nötig}
\end{frame}

\begin{frame}[fragile]
\frametitle{Interface implementieren}
Wie implementiert man ein Interface?
\begin{lstlisting}
    public class Schwert implements Item{
        //inhalt der Klasse
    }
\end{lstlisting}
Nun muss die Klasse Schwert eine Methode \textit{weight()} implementieren, die einen \textit{double} zurück gibt.
\end{frame}

\begin{frame}[fragile]
\frametitle{Interface Klasse A}
\begin{lstlisting}
    class Schwert implements Item{
        String name;
        static final double WEIGHT = 7;
        Schwert(String name){
            this.name = name;
        }
        double weight(){
            return WEIGHT;
        }
    }
\end{lstlisting}
\end{frame}

\begin{frame}[fragile]
\frametitle{Interface Klasse B}
\begin{lstlisting}
    class Heiltrank implements Item {
        String name;
        int count;
        static final double WEIGHT = 2.5;
        Heiltrank(String name, int count){
            this.name = name;
            this.count = count;
        }
        double weight(){
            return count * WEIGHT;
        }
    }
\end{lstlisting}
\end{frame}

\begin{frame}[fragile]
\frametitle{Abstrakte Klassen}
\begin{itemize}
    \item Abstrakte Klassen können nicht instanziiert werden, können aber Methoden implementieren. 
    \item Sie können auch abstrakte Methoden enthalten, die von Unterklassen implementiert werden müssen.
    \item Unterklassen von abstrakten Klassen müssen alle abstrakten Methoden implementieren, es sei denn, sie sind selbst abstrakt.
\end{itemize}
\end{frame}